\documentclass[11pt,a4paper]{article}
\usepackage[utf8]{inputenc}\usepackage[T1]{fontenc}
\usepackage{amsmath,amssymb,bm,array,geometry,color,fancyhdr,hyperref}
\geometry{left=19mm,right=19mm,top=18mm,bottom=19mm}
\definecolor{ink}{RGB}{20,52,72}\definecolor{accent}{RGB}{22,115,139}
\hypersetup{colorlinks=true,urlcolor=accent,linkcolor=ink}
\setlength{\parindent}{0pt}\setlength{\parskip}{3pt}
\setlength{\abovedisplayskip}{6pt}\setlength{\belowdisplayskip}{6pt}\setlength{\jot}{3pt}
\pagestyle{fancy}\fancyhf{}\renewcommand{\headrulewidth}{0pt}
\fancyfoot[L]{\footnotesize pysnspd 1.2.1 / efectivo / \href{Supuestos.pdf}{Supuestos}}\fancyfoot[R]{\small\thepage}
\newcommand{\block}[1]{\par\vspace{4pt}{\large\bfseries\color{ink}#1}\par\vspace{3pt}}
\newcommand{\labelrow}[2]{\(\displaystyle #1\)& #2\\}
\newcommand{\origin}[1]{\par{\footnotesize\textbf{Campo / procedencia:} #1}\par}
\newcommand{\code}[1]{\par{\footnotesize\textbf{Código:} \texttt{#1}}\par}
\newcommand{\dd}{\,\mathrm d}\newcommand{\FD}{f_{\rm FD}}\newcommand{\BE}{n_{\rm BE}}
\newcommand{\Rea}{\operatorname{Re}}\newcommand{\Ima}{\operatorname{Im}}
\newcommand{\refs}{\vfill{\footnotesize\href{Supuestos.pdf}{Hipótesis y límites: Supuestos}; \href{Parametros.pdf}{entradas: Parametros}.}}
\begin{document}
{\LARGE\bfseries\color{ink}Ecuaciones auxiliares durante la evolución}\par
{\small Versión 1.2.1. Funcionales, calor reversible, eventos, potencial, circuito e integrador.}
\renewcommand{\theequation}{2.\arabic{equation}}
\block{Funcional térmico anclado a DC y fuerza variacional}\small
\hypertarget{G-final}{}
\begin{gather}
 U_{\rm BCS}=\sum_i u_i^*\left[
 \frac{\rho_i}{g}-2\sum_kw_k(e_{ik}-\bar\xi_k)+4\sum_kw_ke_{ik}p_{ik}\right],\\
 b_{ij}=d_i-e^{-i\alpha_{ij}}d_j,\quad
 C_{ij}=N_0E_u^2d_f\ell_0^2c_{ij},\quad
 L_i=\tfrac12\sum_{j\sim i}C_{ij}|b_{ij}|^2,\\
 \kappa(\rho,T)=\frac{\pi}{2\sqrt\rho}
       \tanh\frac{\sqrt\rho}{2T/T_c},\quad
 \theta_i=T_{E,i}/T_c,\quad b=T_{\rm ref}/T_c,\quad T_{\rm ref}=T_b,\\
 \mathcal Q_i=\frac{L_i}{u_i^*\rho_i},\quad
 Q_{{\rm dep},i}=E_u\mathcal Q_i
 \quad(\rho_i>0),\\
 \boxed{\mathcal G_i=\kappa(\rho_i,T_{\rm ref})L_i+
 u_i^*[\varphi(\rho_i,\theta_i,\mathcal Q_i)
       -\varphi(\rho_i,b,\mathcal Q_i)]},\quad G=\sum_i\mathcal G_i,\\
 G(T_E=T_{\rm ref})=\sum_i\kappa(\rho_i,T_{\rm ref})L_i,\quad
 \mathcal G_i=\kappa(\rho_i,T_{E,i})L_i+O(L_i^2),\\
 \boxed{U=U_{\rm BCS}+G-\sum_i T_{E,i}G_{T_i}},\qquad
 \boxed{S=\sum_iS_{qp,i}-\sum_iG_{T_i}},\\
 S_{qp,i}=-4N_0E_uV_ik_B\sum_kw_k[p_{ik}\log p_{ik}
                    +(1-p_{ik})\log(1-p_{ik})],\\
 K^L_{ij}=\tfrac12(G_{L_i}+G_{L_j}),\qquad
 G_{L_i}=\partial_{L_i}\mathcal G_i,\quad G_{\rho_i}=\partial_{\rho_i}\mathcal G_i,\\
 \boxed{\mathcal F_i=u_i^*\left[\frac1g-\sum_k
 \frac{w_k(1-2p_{ik})}{e_{ik}}\right]d_i
 +G_{\rho_i}d_i+\sum_{j\sim i}C_{ij}K^L_{ij}b_{ij}},\\
 I^s_{ij}=-\frac{2e}{\hbar}\partial_{\alpha_{ij}}G
 =\frac{4e}{\hbar}C_{ij}K^L_{ij}
        \Ima(d_i^*e^{-i\alpha_{ij}}d_j),\\
 \dd U=2\Rea\sum_i\mathcal F_i^*\dd d_i+
      \sum_{ik}\epsilon_{ik}\dd p_{ik}-\sum_iT_{E,i}\dd G_{T_i}
      +\sum_{\langle ij\rangle}G_{\alpha_{ij}}\dd\alpha_{ij}.
\end{gather}
\par\begin{tabular}{@{}p{.29\linewidth}p{.67\linewidth}@{}}
\hline
\labelrow{g,\mathcal G_i,G_{\rho_i},G_{L_i}}{emparejamiento; energía libre nodal [J]; derivadas [J], [1]}
\labelrow{C_{ij},c_{ij},L_i}{coeficiente [J], conductancia geométrica, peso nodal [J]}
\labelrow{T_{\rm ref},\varphi,\mathcal Q_i}{referencia DC congelada [K]; corrección y depareo adimensionales}
\labelrow{\alpha_{ij}}{\((2e/\hbar)\int_i^j\mathbf A\cdot\dd\mathbf l\); cero en los casos actuales}
\hline
\end{tabular}\par
\origin{BCS/KWT efectivo. El matching propio \(H_i^{\mathrm{match}}=\kappa_bL_i-u_i^*\varphi_b\ge0\) conserva DC, pero \(\mathcal G_i\to H_i^{\mathrm{match}}\) a \(T_E\to\infty\): rigidez residual no pequeña en hotspots. Hipótesis exploratoria y límites: \href{Supuestos.pdf\#nameddest=S15}{S15}; no Usadel exacto ni admisión cuantitativa de latencia.}
\refs\newpage\normalsize
\hypertarget{matsu-runtime}{}\block{Corrección de depareo: cálculo durante cada evaluación}\small
\begin{gather}
 a=\sqrt\rho,\quad\theta=T/T_c,\quad
 \varpi_\nu=(2\nu+1)\pi\theta,\quad \nu=0,1,\ldots,\\
 f_\nu(\chi;\mathcal Q)=-2a\sin\chi+2\varpi_\nu(1-\cos\chi)
                                      +\mathcal Q\sin^2\chi,\\
 \mathcal W(\rho,\theta,\mathcal Q;\boldsymbol\chi)
 =\rho\log\theta+2\pi\theta\sum_{\nu\ge0}
           [\rho/\varpi_\nu+f_\nu(\chi_\nu;\mathcal Q)],\\
 \partial_{\chi_\nu}\mathcal W=0
 \ \Longleftrightarrow\
 \varpi_\nu\sin\chi_\nu+\mathcal Q\sin\chi_\nu\cos\chi_\nu-a\cos\chi_\nu=0,\\
 v_\nu=\tan\chi_\nu/a,\quad
 \boxed{\varpi_\nu v_\nu+\frac{\mathcal Qv_\nu}{\sqrt{1+\rho v_\nu^2}}=1},\quad
 \frac1{\varpi_\nu+\mathcal Q}\le v_\nu\le\frac1{\varpi_\nu},\\
 c_\nu=(1+\rho v_\nu^2)^{-1/2},\quad r_\nu=v_\nu c_\nu,\quad
 H_\nu=\varpi_\nu+\mathcal Qc_\nu^3>0,\quad
 E_\nu^M=\sqrt{\varpi_\nu^2+\rho},\\
 r_{\nu0}=1/E_\nu^M,\quad c_{\nu0}=\varpi_\nu/E_\nu^M,\quad
 B_\nu=E_\nu^M\rho(r_\nu-r_{\nu0})^2
 \left[1+\frac{\rho(r_\nu+r_{\nu0})^2}{(c_\nu+c_{\nu0})^2}\right]
 +\mathcal Q\rho r_\nu^2,\\
 \boxed{\varphi(\rho,\theta,\mathcal Q)=
 \mathcal W|_{\boldsymbol\chi_{\mathcal Q}}-\mathcal W|_{\boldsymbol\chi_0}
 =2\pi\theta\sum_{\nu\ge0}B_\nu},\quad
 \varphi(\rho,\theta,0)=0,\\
 \varphi_\rho=2\pi\theta\sum_\nu(r_{\nu0}-r_\nu),\quad
 \varphi_{\mathcal Q}=2\pi\theta\sum_\nu\rho r_\nu^2,\\
 \partial_\rho v_\nu=\frac{\mathcal Qv_\nu^3c_\nu^3}{2H_\nu},\quad
 \partial_{\mathcal Q}v_\nu=-v_\nu c_\nu/H_\nu,\quad
 \partial_{\varpi_\nu}v_\nu=-v_\nu/H_\nu,\\
 \mathbf D_{\le2}\varphi=
 (\varphi,\varphi_\rho,\varphi_{\mathcal Q},\varphi_\theta,
 \varphi_{\rho\rho},\varphi_{\rho\mathcal Q},\varphi_{\mathcal Q\mathcal Q},
 \varphi_{\rho\theta},\varphi_{\mathcal Q\theta},\varphi_{\theta\theta}),\\
 \varphi_{\theta\theta}=\frac{2\pi}{\theta}\sum_\nu
 \left[4\varpi_\nu(c_{\nu0}-c_\nu)
 -2\varpi_\nu^2\rho\left(\frac{r_\nu^2c_\nu}{H_\nu}-r_{\nu0}^3\right)\right].
\end{gather}
{\small \(\rho\) permanece independiente: se estacionariza el ángulo espectral, no el gap KWT.
Todas las derivadas de \(\mathbf D_{\le2}\varphi\) se obtienen analíticamente con la raíz implícita. La suma depende del estado: no es una tabla precomputada.}
\begin{gather}
 N_M=\max\!\left(64,\left\lceil\frac{32\sqrt\rho}{2\pi\theta}\right\rceil\right),
 \quad q_M=\mathcal Q/(2\pi\theta),\\
 \varphi=\rho[\psi(\tfrac12+q_M)-\psi(\tfrac12)]
 +2\pi\theta\!\sum_{\nu<N_M}\!
 [B_\nu-\rho(\varpi_\nu^{-1}-(\varpi_\nu+\mathcal Q)^{-1})]+\mathcal T_5,\\
 \mathcal T_5=2\pi\theta\sum_{\nu\ge N_M}\sum_{m=2}^{5}
 \rho^m[\mathcal B_m(\varpi_\nu,\mathcal Q)-\mathcal B_m(\varpi_\nu,0)],\\
 \mathcal B_m(\omega,Q):=
 [x^m]\,\min_\chi\{-2\sqrt x\sin\chi+
                  2\omega(1-\cos\chi)+Q\sin^2\chi\},\\
 2\pi\theta\sum_{\nu\ge N_M}(\varpi_\nu+\mathcal Q)^{-p}
 =(2\pi\theta)^{1-p}\zeta(p,N_M+\tfrac12+q_M).
\end{gather}
{\small \(\psi,\zeta\): digamma y zeta de Hurwitz; \([x^m]\): coeficiente de la serie en \(x\). Cola hasta \(\rho^5\), Euler--Maclaurin de seis términos; sin expansión en \(\mathcal Q/\varpi\).}
\origin{\href{https://arxiv.org/html/1207.6421}{Clem--Kogan (2012)}, ecs.~7--14 y apéndice D: forma variacional reconstruida del problema espectral. No se usa la energía ya auto-consistente de su ec.~116 como potencial fuera de equilibrio.}
\code{finite\_depairing\_sums.py}
\refs\newpage\normalsize
\block{Derivadas del cierre y calor reversible}\small
{\small Nodo fijo; \(u=u_i^*,\,\rho=\rho_i,\,L=L_i,\,T=T_{E,i}\),
 \(\Delta\varphi=\varphi(\rho,T/T_c,\mathcal Q)-\varphi(\rho,b,\mathcal Q)\),
 \(\mathcal Q=L/(u\rho)\); \(\kappa_b=\kappa(\rho,T_{\rm ref})\).
 Subíndices: derivadas parciales a las otras variables explícitas fijas.}
\begin{gather}
 G_L=\kappa_b+\Delta\varphi_{\mathcal Q}/\rho,\quad
 G_\rho=L\kappa_{b,\rho}+u(\Delta\varphi_\rho-\mathcal Q\Delta\varphi_{\mathcal Q}/\rho),\\
 G_{LL}=\Delta\varphi_{\mathcal Q\mathcal Q}/(u\rho^2),\quad
 G_{\rho L}=\kappa_{b,\rho}+\Delta\varphi_{\rho\mathcal Q}/\rho
       -(\Delta\varphi_{\mathcal Q}+\mathcal Q\Delta\varphi_{\mathcal Q\mathcal Q})/\rho^2,\\
 G_{\rho\rho}=L\kappa_{b,\rho\rho}+u\!\left[
 \Delta\varphi_{\rho\rho}-\frac{2\mathcal Q}{\rho}\Delta\varphi_{\rho\mathcal Q}
 +\frac{\mathcal Q^2\Delta\varphi_{\mathcal Q\mathcal Q}
              +2\mathcal Q\Delta\varphi_{\mathcal Q}}{\rho^2}\right],\\
 G_T=\frac{u}{T_c}\varphi_\theta,\quad G_{TT}=\frac{u}{T_c^2}\varphi_{\theta\theta},\quad
 G_{LT}=\frac{\varphi_{\mathcal Q\theta}}{\rho T_c},\quad
 G_{\rho T}=\frac{u}{T_c}\left(\varphi_{\rho\theta}
                    -\frac{\mathcal Q}{\rho}\varphi_{\mathcal Q\theta}\right),\\
 f_{ik}=\FD(E_{ik},T_i),\quad f_{e,ik}=-\frac{E_u}{k_BT_i}f_{ik}(1-f_{ik}),\\
 C_{{\rm FD},i}=\sum_k\epsilon_{ik}\frac{E_{ik}}{k_BT_i^2}f_{ik}(1-f_{ik}),\quad
 B_i=4u_i^*\sum_k\frac{w_k}{2e_{ik}}[p_{ik}-f_{ik}-e_{ik}f_{e,ik}],\\
 C_{{\rm FD},i}\dot T_i=\sum_k\epsilon_{ik}\dot p_{ik}+B_i\dot\rho_i,\quad
 \dot\rho_i=2\Rea(d_i^*\dot d_i),\\
 \dot L_i=\sum_{j\sim i}C_{ij}\Rea[
 b_{ij}^*(\dot d_i-e^{-i\alpha_{ij}}\dot d_j)]\quad(\dot\alpha=0),\\
 A_i^{G}=G_{L_iT_i}\dot L_i+G_{\rho_iT_i}\dot\rho_i,\quad
 \boxed{\dot G_{T_i}=A_i^{G}+G_{T_iT_i}\dot T_i},\\
 P_i^{bare}=\sum_k\epsilon_{ik}
        (C^{ep}_{ik}+C^{tr}_{ik}+C^{BGK}_{ik}),\\
 \boxed{C_{{\rm eff},i}\dot T_i
 =P_i^{bare}+Q_{\Delta,i}+P_{J,i}+B_i\dot\rho_i+T_iA_i^G},\\
 \boxed{C_{{\rm eff},i}=C_{{\rm FD},i}-T_iG_{T_iT_i}>0},\quad
 \boxed{G_{L_i}>0,\quad G_{L_i}+2L_iG_{L_iL_i}>0}.
\end{gather}
\block{Límites normales, sin piso de gap}
\begin{gather}
 \rho=0,\ L>0:\quad
 G=L\kappa(0,T_{\rm ref}),\quad
 G_\rho=L\kappa_\rho(0,T_{\rm ref})+u\log(T_{\rm ref}/T),\\
 G_T=G_{TT}=G_{LT}=0,\quad G_{\rho T}=-u/T,\quad
 \rho=L=0:\quad \mathcal F^{grad}=0,\ I^s=0.
\end{gather}
{\small El nodo normal aislado tiene curvatura finita. En el origen homogéneo la curvatura depende de la dirección; el Hessiano de gradiente pequeño sólo precondiciona. \(\dot T_E\) es algebraica, no un estado independiente. Rechazar capacidad o rigidez principal no positivas; no recortar \(T_E\), \(G_L\) ni \(C_{\rm eff}\).}
\par\begin{tabular}{@{}p{.29\linewidth}p{.67\linewidth}@{}}
\hline
\labelrow{C_{{\rm FD},i},C_{{\rm eff},i}}{capacidades del momento FD y efectiva [J/K]}
\labelrow{B_i,A_i^G,\dot L_i}{coeficientes [J], [W/K], tasa geométrica [J/s]}
\labelrow{G_{\rho\rho},G_{\rho L},G_{LL}}{derivadas [J], [1], [J\(^{-1}\)]}
\labelrow{G_T,G_{TT},G_{LT},G_{\rho T}}{derivadas [J/K], [J/K\(^{2}\)], [K\(^{-1}\)], [J/K]}
\hline
\end{tabular}\par
\origin{Regla de la cadena de \(G(\rho,L,T_E)\), momento energético nativo y guardia de ellipticidad del cierre efectivo.}
\refs\newpage\normalsize
\hypertarget{carrier-exacto}{}\block{Portador de calor con actividades de Pauli}
{\small Nodo fijo; $T=T_E$, $f_k=\FD(E_k,T)$. $P=Q_\Delta+P_J+T\dot G_T$ [W]. El signo selecciona una actividad de nacimiento o de muerte.}
\begin{gather}
 a_k=f_k(1-p_k),\qquad b_k=p_k(1-f_k),\\
 \boxed{h_k^+=\frac{E_ka_k}{\sum_l\epsilon_lE_la_l}},\qquad
 \boxed{h_k^-=\frac{E_kb_k}{\sum_l\epsilon_lE_lb_l}},\\
 \boxed{H_k=P h_k^{\operatorname{sgn}P}},\qquad
 \sum_k\epsilon_k h_k^\pm=1,\qquad
 h_k^+\big|_{p_k=1}=h_k^-\big|_{p_k=0}=0,\\
 p_k=f_k\ \Longrightarrow\
 h_k^\pm=\frac{\partial_T f_k}{C_{\rm FD}},\quad
 \partial_T f_k=\frac{E_k}{k_BT^2}f_k(1-f_k),\quad
 C_{\rm FD}=\sum_k\epsilon_k\partial_Tf_k .
\end{gather}
{\small Se exige $T>0$, $0\le p\le1$ y normalizador finito positivo. No se evalúan logaritmos para construir $H$; no hay clipping ni piso de ocupación. $P=0$ da $H=0$ y usa la rama $+$.}
\block{Desigualdad entrópica y energía}
\begin{gather}
 \sigma_k=\frac{\epsilon_k k_B}{E_k}\log\frac{1-p_k}{p_k},\quad
 \log(a_k/b_k)=\log\frac{1-p_k}{p_k}-\frac{E_k}{k_BT},\\
 \sum_k\epsilon_k(a_k-b_k)=\sum_k\epsilon_k(f_k-p_k)=0,\\
 a_k\log(a_k/b_k)\ge a_k-b_k,\qquad
 b_k\log(a_k/b_k)\le a_k-b_k,\\
 \boxed{\dot S_{\rm exc}=P\left(\sum_k\sigma_kh_k^{\operatorname{sgn}P}-T^{-1}\right)\ge0},\\
 \dot S_{\rm source}-\dot G_T
   =\frac{Q_\Delta+P_J}{T}+\dot S_{\rm exc}\ge0,\\
 \sum_k\epsilon_{ik}C^{BGK}_{ik}=0,\qquad
 \sum_{ik}\epsilon_{ik}C^{tr}_{ik}=0,\\
 \sum_k\epsilon_{ik}C^{ep}_{ik}
       +u_i^*\sum_qc_q\omega_qC^{ph,ep}_{iq}=0,\\
 2\Rea\sum_i\mathcal F_i^*(D_td_i)=-\sum_iQ_{\Delta,i},\quad
 \dot U+\dot U_{ph}=I_sV_{\rm dev}-P_{\rm esc}-P_{\rm cuts}.
\end{gather}
{\small La prueba entrópica se formula en el interior; en los extremos se toma el límite unilateral. El RHS sí se evalúa en los extremos. La igualdad entrópica se recupera en FD, no se impone fuera de equilibrio.}
\par\begin{tabular}{@{}p{.29\linewidth}p{.67\linewidth}@{}}
\hline
\labelrow{a_k,b_k}{actividades de aporte y retiro, adimensionales}
\labelrow{h_k^\pm,H_k}{portador [J$^{-1}$], fuente de ocupación [s$^{-1}$]}
\labelrow{\epsilon_k,\sigma_k,\dot S_{\rm exc}}{energía [J], momento entrópico [J/K], producción [W/K]}
\hline
\end{tabular}\par
\origin{Exclusión de Pauli, momento energético FD y desigualdad de Gibbs. Reparto espectral constitutivo propio; no es una colisión microscópica derivada. Sustituye el portador 1.1.0, que no preservaba soporte al enfriar.}
\code{thermal\_stiffness.py: entropy\_compatible\_carrier}
\refs\newpage
\block{Eventos electrón--fonón compartidos}
\begin{gather}
 a=S:(k,l),\ k<l;\quad a=R:(k,l),\ k\leq l,\\
 \omega^S=e_{il}-e_{ik},\quad \omega^R=e_{il}+e_{ik},\quad
 M^{S/R}=1\mp\rho_i/(e_{ik}e_{il}),\\
 \widetilde w^S=w_kw_l,\quad \widetilde w^R=w_kw_l/(1+\delta_{kl}),\\
 {\cal A}^{+}=\prod_q(1+n_{iq})^{\nu^{ph}_{qa}},\quad
 {\cal A}^{-}=\prod_qn_{iq}^{\nu^{ph}_{qa}},\\
 b^S=p_{il}(1-p_{ik}){\cal A}^+-p_{ik}(1-p_{il}){\cal A}^-,\\
 b^R=p_{ik}p_{il}{\cal A}^+-(1-p_{ik})(1-p_{il}){\cal A}^-,\\
 r^{S/R}_{ia}=\frac{8\pi E_u}{\hbar}\widetilde w^{S/R}
                  \alpha^2F(E_u\omega^{S/R})M^{S/R}b^{S/R},\\
 \nu^e_{ka}=+1\ (S,\text{inferior}),\ -1\ (S,\text{superior}),\
 -1\ (R,\text{cada participante}),\\
 \nu^e_{ka}=-2\quad(R,k=l),\quad
 \beta=\frac{\omega_{ia}-\omega_q}{\omega_{q+1}-\omega_q},\\
 \nu^{ph}_{qa}=1-\beta,\quad\nu^{ph}_{q+1,a}=\beta,\quad
 \sum_q\omega_q\nu^{ph}_{qa}=\omega_{ia}.
\end{gather}
\block{Transporte electrónico a energía común}
\begin{gather}
 e_{\ell,-}=\max(e_{i0},e_{j0}),\quad
 e_{\ell,+}=\min(e_{i,K-1},e_{j,K-1}),\\
 e_{\ell a}=\tfrac12(e_{\ell,-}+e_{\ell,+})
             +\tfrac12(e_{\ell,+}-e_{\ell,-})x_a,\quad
 w_{\ell a}=\tfrac12(e_{\ell,+}-e_{\ell,-})w_a^{GL},\\
 e_{\ell a}=\sum_k\beta^{(i)}_{ka}e_{ik}
            =\sum_k\beta^{(j)}_{ka}e_{jk},\quad\sum_k\beta^{(i)}_{ka}=1,\\
 {\cal P}_{ia}=\prod_kp_{ik}^{\beta^{(i)}_{ka}},\
 {\cal Q}_{ia}=\prod_k(1-p_{ik})^{\beta^{(i)}_{ka}},\\
 J_{\ell a}=4D\frac{A_\ell^{face}}{L_\ell}w_{\ell a}
             ({\cal P}_{ia}{\cal Q}_{ja}-{\cal P}_{ja}{\cal Q}_{ia}).
\end{gather}
\par\begin{tabular}{@{}p{.29\linewidth}p{.67\linewidth}@{}}
\hline
\labelrow{S,R,M,\widetilde w}{dispersión, recombinación, coherencia BCS, peso de par}
\labelrow{\alpha^2F,\omega,\nu}{Eliashberg adimensional, energía/\(E_u\), estequiometría}
\labelrow{A_\ell^{face},L_\ell,D}{área [m\(^{2}\)], distancia [m], difusión [m\(^{2}\)/s]}
\labelrow{x_a,w_a^{GL},\beta}{cuadratura de referencia e interpolación energética}
\hline
\end{tabular}\par
\origin{\href{https://arxiv.org/pdf/1611.06060}{Vodolazov (2017)}, cinética BCS, y volúmenes finitos propios. No hay difusión lateral fonónica en este cierre.}
\refs\newpage
\block{Potencial cuasineutro y cortes superconductores}
\begin{gather}
 G_\ell^{n}=c_\ell/R_\square,\quad
 I_\ell^n=G_\ell^{n}(\phi_i-\phi_j),\quad
 I_\ell=I_\ell^s+I_\ell^n,\\
 \phi=P\Phi,\quad B_g=P^{\mathsf T}B,\quad
 (B_gG^nB_g^{\mathsf T})\Phi=\mathbf s_gI_s-B_gI^s,\\
 s_{g,L}=1,\quad s_{g,R}=-1,\quad s_{g,\mathrm{int}}=0,\quad\phi_R=0,\\
 P_{J,\ell}=G_\ell^n(\phi_i-\phi_j)^2,\quad
 P_{J,i}=\tfrac12\sum_{\ell\ni i}P_{J,\ell},\\
 I_sV_{\rm dev}=\sum_\ell I_\ell^s(\phi_i-\phi_j)+\sum_\ell P_{J,\ell},
 \quad V_{\rm dev}=\phi_L-\phi_R,\\
 d_i=|d_{{\rm DC},i}|e^{i\theta_{L/R}},\quad
 \dot\theta_{L/R}=-2e\phi_{L/R}/\hbar,\quad D_td_i=0\quad(i\in L/R),\\
 p_{ik}=\FD(E_{ik},T_b),\ n_{iq}=\BE(\Omega_q,T_b),\
 \dot p_{ik}=\dot n_{iq}=\dot T_{E,i}=0\quad(i\in L/R),\\
 \hat{\mathbf n}\cdot\mathbf j_{\rm tot}
 =\hat{\mathbf n}\cdot\boldsymbol\Phi_p=0\quad(\text{bordes laterales}),\\
 P_{\rm cuts}=\sum_{i\in L\cup R}\left[
       \sum_k\epsilon_{ik}\dot p_{ik}^{raw}
       +u_i^*\sum_qc_q\omega_q\dot n_{iq}^{raw}\right],\\
 P_{\rm esc}=\sum_{iq}\frac{u_i^*c_q\omega_q}{\tau_{{\rm esc},i}}
             [n_{iq}-\BE(\Omega_q,T_b)],\\
 \dot p^{raw}=C^{ep}+C^{tr}+C^{BGK}
       +h(Q_\Delta+P_J+T_b\dot G_T)\quad(\text{cortes}).
\end{gather}
{\small \(\dot G_T\) puede ser distinto de cero en un corte de temperatura fija: cambia su enlace con el interior. Se incluye en el RHS bruto antes de congelar \(p,n\). No se añade otra vez a \(P_{\rm cuts}\).}
\par\begin{tabular}{@{}p{.29\linewidth}p{.67\linewidth}@{}}
\hline
\labelrow{B,B_g,P,\Phi}{incidencias, prolongación equipotencial, potencial reducido [V]}
\labelrow{\mathbf s_g,G_\ell^n}{inyección de corriente, conductancia normal fija [S]}
\labelrow{L,R,\theta_{L/R}}{cortes completos, fases Josephson}
\labelrow{\boldsymbol\Phi_p,P_{\rm cuts},P_{\rm esc}}{flujo de población, potencias salientes [W]}
\hline
\end{tabular}\par
\origin{Continuidad cuasineutra, red óhmica, identidad gauge de \(G\). Continuación superconductora rígida, no contacto metálico; véase \href{Supuestos.pdf}{Supuestos}.}
\refs\newpage
\block{Circuito de tres estados y observables}
\begin{gather}
 I_{\rm RF}=I_b-I_s,\quad V_{\rm out}=R_l(I_b-I_s),\\
 L_b\dot I_b=V_b-R_bI_b-v_c-R_l(I_b-I_s),\\
 L_{\rm ext}\dot I_s=v_c+R_l(I_b-I_s)-V_{\rm dev},\quad
 C_c\dot v_c=I_b-I_s,\\
 U_c=\tfrac12L_bI_b^2+\tfrac12L_{\rm ext}I_s^2+\tfrac12C_cv_c^2,\\
 \dot U_c=V_bI_b-R_bI_b^2-R_l(I_b-I_s)^2-I_sV_{\rm dev},\\
 U_{\rm tot}=U+U_{ph}+U_c,\quad
 U_{ph}=\sum_{iq}u_i^*c_q\omega_qn_{iq},\\
 P_{\rm ext}=V_bI_b-R_bI_b^2-R_l(I_b-I_s)^2-P_{\rm esc}-P_{\rm cuts},\\
 {\cal R}_P=\dot U_{\rm tot}-P_{\rm ext},\quad
 {\cal R}_U=U_{\rm tot}(t)-U_{\rm tot}(0)-\int_0^tP_{\rm ext}\dd t .
\end{gather}
\block{Cruce del puerto, confirmación y ventana extra}
\begin{gather}
 V_{\rm sig}=V_{\rm out}-V_{\rm DC},\quad V_{\rm th}=0.1\,\mathrm{mV},\\
 V_k<V_{\rm th}\leq V_{k+1},\quad
 t_\times=t_k+\frac{V_{\rm th}-V_k}{V_{k+1}-V_k}(t_{k+1}-t_k),\\
 V_{\rm sig}\geq0.95V_{\rm th}\quad
    [t_\times,t_\times+0.5\,\mathrm{ps}],\quad
 t_{\rm end}\geq t_\times+5\,\mathrm{ps},\\
 t_0=0\ (\text{postcascada}),\quad\tau_{\rm lat}=t_\times-t_0,\quad
 \Delta\tau=\tau_{\rm lat}(1550)-\tau_{\rm lat}(775),\\
 I_b>0,\ I_s>0,\qquad t_{\max}^{\rm sin\ cruce}=100\,\mathrm{ps}.
\end{gather}
{\small \(V_{\rm out}\) no incluye amplificador, filtro ni ruido; no equivale al voltaje de osciloscopio de Korzh. No cruce al techo: censura. Fallo de capacidad/integración: fallo del candidato, no censura física. Sin reset eléctrico completo.}
\par\begin{tabular}{@{}p{.29\linewidth}p{.67\linewidth}@{}}
\hline
\labelrow{I_b,I_s,I_{\rm RF}}{corrientes de bias, detector y lectura [A]}
\labelrow{V_b,v_c,V_{\rm dev},V_{\rm out}}{voltajes de fuente, capacitor, región 2D y carga [V]}
\labelrow{R_b,R_l,L_b,L_{\rm ext},C_c}{resistencias [\(\Omega\)], inductancias [H], capacitor [F]}
\labelrow{{\cal R}_P,{\cal R}_U,t_\times}{defectos independientes [W], [J], cruce [s]}
\hline
\end{tabular}\par
\origin{Kirchhoff y energía electromagnética; circuito de la memoria, parámetros heredados explícitos en \href{Parametros.pdf}{Parametros}, no setup completo publicado de Korzh.}
\refs\newpage
\block{Integración conjunta y admisión numérica}\small
\begin{gather}
 {\cal X}=(d,p,n,I_b,I_s,v_c),\quad
 \delta{\cal X}=h{\cal F}_h({\cal X}^{mid}),\quad
 {\cal X}^{mid}=\tfrac12({\cal X}^{m+1}+{\cal X}^m),\\
 d_{L/R}^{m+1}=d_{L/R}^me^{i\delta\theta},\quad
 d_{L/R}^{mid}=d_{L/R}^me^{i\delta\theta/2},\\
 \gamma_h(\tau)=\frac2h\tanh\frac{h}{2\tau},\quad
 r_a^h=\lambda_ar_a,\quad J_a^h=\lambda_a^{tr}J_a,\qquad
 \lambda_a=\min_{\mu\in{\cal C}(a)}\min(1,0.95\,b_\mu/L_\mu),\\
 g_a=h{\cal F}_{h,a}-\delta a,\quad
 {\cal E}_p=\sum_{ik}\epsilon_{ik}^{mid}|g_{p,ik}|,\quad
 {\cal E}_{ph}=\sum_{iq}u_i^*c_q\omega_q|g_{n,iq}|,\\
 E_{\rm ref}=\begin{cases}\eta E_\gamma,&\text{fotón},\\
                         \sum_i u_i^*,&\text{DC},\end{cases}\\
 \max(\|g_d\|_\infty,\|g_c\|_\infty,
       {\cal E}_p/E_{\rm ref},{\cal E}_{ph}/E_{\rm ref},r_{p\to d}^{full})
          <\epsilon_{\rm NL}(h),\\
 \delta W_m=h_mP_{\rm ext}^{mid},\quad
 \varepsilon_{E,m}=U_{\rm tot}^{m+1}-U_{\rm tot}^{m}-\delta W_m,\\
 |\varepsilon_{E,m}|\leq\epsilon_{\rm guard}
 (E_{\rm ref}+\sum_{\ell<m}|\delta W_\ell|+|\delta W_m|),\\
 0\leq p^{m+1}\leq1,\quad n^{m+1}\geq0,\quad
 C_{\rm eff}({\cal X}^{mid}),C_{\rm eff}({\cal X}^{m+1})>0,\quad G_L>0,\quad G_L+2LG_{LL}>0,\\
 \epsilon_{\rm NL}(h)=2\,10^{-7}\min(1,h/(0.02\,{\rm ps})),\quad
 \epsilon_{\rm guard}=2.5\,10^{-6},\quad h_{\max}=0.02\,\mathrm{ps},\\
 \delta p_{ik}=\tfrac12g_{p,ik},\qquad
 \delta T_i=\frac{\sum_k\epsilon_{ik}\delta p_{ik}}{C_{{\rm FD},i}},\\
 \delta\mathcal F=(D_p\mathcal F)_{\mathbf T}\delta p+
 (D_{\mathbf T}\mathcal F)_p\delta\mathbf T,\qquad
 \delta I^s=D_{\mathbf T}I^s\,\delta\mathbf T,\\
 (B_gG^nB_g^{\mathsf T})\delta\Phi=-B_g\delta I^s,\qquad
 \delta\phi=P\delta\Phi,\\
 \delta R_{d,i}=-\mathcal M_i\delta\mathcal F_i
 -(\delta\mathcal M_i)\mathcal F_i
 -i\frac{2e}{\hbar}\delta\phi_i d_i\quad(i\ {\rm interior}),\\
 \delta v_i=\frac{\max(T_i+\delta T_i,T_b)-\max(T_i,T_b)}{T_c},
 \qquad \delta R_{\theta,c}=-\frac{2e}{\hbar}\delta\phi_c,\\
 \boxed{r_{p\to d}^{full}:=h\max\!
 \left(\|\delta R_{d,\rm int}\|_\infty,
       \|\delta R_{\theta,\rm cuts}\|_\infty\right)
 <\epsilon_{\rm NL}(h)}.
\end{gather}
{\small $h$ [s]: paso; $h_{ik}$ [J$^{-1}$]: portador; $g_d,g_c$: residuos escalados.
$b_\mu,L_\mu$: presupuesto/pérdida de tasa de estados consumidos $\mathcal C(a)$; $\lambda$ conserva su estequiometría.
La fuente electrónica firmada queda fuera del limitador. Hessiano a $T_E$ fijo: sólo precondicionador.}
{\small $D_p,D_{\mathbf T}$: derivadas analíticas; $(D_p\mathcal F)_{\mathbf T}\delta p=2u^*d\sum_kw_k\delta p_k/e_k$.
$\mathcal M$ [J$^{-1}$s$^{-1}$]: operador KWT; $\delta\mathcal M$ usa el incremento $\delta v$ definido arriba.
$\delta\mathcal F,\delta I^s,\delta\phi,\delta R_d,\delta R_\theta$ tienen unidades J, A, V, s$^{-1}$, s$^{-1}$.
$\delta p=g_p/2$ es dirección: no se evalúa $p+\delta p$. El indicador no acota el error temporal.}
{\small Exigir al menos una corrección conjunta y el RHS completo convergido.
La búsqueda interna mantiene soporte, $C_{\rm eff}>0$ y ambas rigideces positivas, sin pisos.
Si $\Delta t_{\rm adapt}<10^{-6}$ ps: checkpoint y parada numérica; nunca aumentar el paso ni asignar detección.}
\begin{gather}
 x_k=(\delta d_k,\delta p_k,\delta n_k/(1+n^m),\delta c_k),\\
 \gamma_k=\arg\min_\gamma\|\Delta R_k\gamma-r_k\|_2,\qquad
 \delta x_k=r_k-(\Delta X_k+\Delta R_k)\gamma_k,\quad m_{\rm AA}\le10.
\end{gather}
{\small \(r_k\): mapa conjunto precondicionado de gap/circuito y residuos de poblaciones;
\(\Delta X_k,\Delta R_k\): diferencias de iteraciones. Anderson acelera las incógnitas conjuntas, no modifica el RHS, tolerancias ni ecuación aceptada. Soporte y guardias se comprueban antes de aplicar la corrección.}
\block{Comparación mínima antes de admitir latencias}
\[
 h_{\max}\to h_{\max}/2,\quad
 h_{\rm mesh}:5\to3.5\,\mathrm{nm},\quad
 (K,Q,N_{\rm face})\ \text{refinados};\qquad
 \{\tau_{\rm lat},V_{\rm out},\Delta_{\min},C_{\rm eff},{\cal R}_U\}.
\]
{\small El cierre cuadrático térmico anterior desarrolló una frontera de capacidad en los anclajes. El modelo 1.2.1 cambia el funcional mediante el ajuste DC declarado; no se admiten trayectorias históricas bajo la energía nueva. Validación fotónica y convergencia de observables pendientes.}
\origin{Punto medio implícito, Newton disperso y limitación de producción/destrucción. Conservación continua no garantiza convergencia temporal ni validez de la hipótesis mesoscópica.}
\code{general\_thermal\_gradient.py; resummed\_thermal.py; effective\_coupled.py}
\refs

\clearpage
\hypertarget{dimensional-1d}{}\hypertarget{hybrid-1d}{}
{\LARGE\bfseries\color{ink}Relevo conservativo y condición inicial}\par
\renewcommand{\theequation}{D2.\arabic{equation}}\renewcommand{\theHequation}{D2.\arabic{equation}}\setcounter{equation}{0}
\block{Remapeo positivo por volúmenes ordenados en \(x\)}
\begin{gather}
 a_i=\sum_{r\leq i}V_r^{(2)},\quad b_j=\sum_{r\leq j}V_{r,C},\\
 B_{ji}=\big|[a_{i-1},a_i]\cap[b_{j-1},b_j]\big|,\quad {\cal R}_{ji}=B_{ji}/V_j^{(1)},\quad V_{j,C}=Wd_f|{\rm celda}_j\cap C|,\\
 {\cal R}_{ji}\geq0,\quad \sum_i{\cal R}_{ji}=f_j=V_{j,C}/V_j^{(1)},\quad
 \sum_jV_j^{(1)}{\cal R}_{ji}=V_i^{(2)},\\
 d_j^0=d_j^{{\rm DC},1}
     \sum_i{\cal R}_{ji}{d_i^{(2)}\over d_i^{{\rm DC},2}}+(1-f_j)d_j^{\rm ext},\quad
 p_{jk}^0=\sum_i{\cal R}_{ji}p_{ik}^{(2)}+(1-f_j)p_{jk}^{\rm ext},\\
 n_{ja}^0=\sum_i{\cal R}_{ji}n_{ia}^{(2)}+(1-f_j)n_{ja}^{\rm ext},\\
 (I_b,I_s,v_c)_1=(I_b,I_s,v_c)_2.
\end{gather}
{\small Remapeo longitudinal de volumen, no promedio puntual exacto en malla
irregular. Se heredan fases relativas a DC y poblaciones rígidas de los cortes.
El núcleo \(C=[350,650]\) nm se remapea por volumen; las colas se suponen frías, en la rama estable al \(I_s\) de transferencia. El perfil relativo a DC conserva fase y evita cambiar el núcleo a FD/BE. Ext significa cola fría al iniciar la réplica y su estado nativo evolucionado al relevo.}
\[
 p^{\rm cold}=\FD(d^{\rm cold},T_b),\quad n^{\rm cold}=\BE(T_b),\quad
 \delta q=q_{\rm grid}(I_s)-q_{\rm grid}(I_{\rm DC}).
\]
{\small Fase de cada cola: fase DC más el incremento del corte central y
\(\delta q(x-x_{\rm interfaz})\); módulo DC multiplicado por
\(a_{\rm grid}(I_s)/a_{\rm grid}(I_{\rm DC})\). Cortes exteriores rígidos
conservan módulo y poblaciones de su DC. La memoria de las colas no se conoce.}
\block{Reconciliación energética sólo durante el cambio de malla}
\begin{gather}
 {\cal E}_q=U_{{\rm BCS},q}+G_q-\sum_jT_{E,j}G_{T,j}+U_{{\rm ph},q},\\
 \delta{\cal E}_q={\cal E}_q(X)-{\cal E}_q(X_q^{\rm DC}),\quad
 \delta{\cal E}_1(X_1^+) =\delta{\cal E}_2(X_2^-)+\tfrac12\Delta L(I_s^2-I_{\rm DC}^2),\\
 \Delta L=L'_{\rm DC}(L_1-L_2),\quad
 L_{{\rm ext},1}=L_{{\rm ext},2}-\Delta L,\\
 {\rm logit}\,p_{jk}^{+}={\rm logit}\,p_{jk}^{0}+\lambda e_{jk},\quad
 \log(1+1/n_{ja}^{+})=\log(1+1/n_{ja}^{0})-\lambda\bar\Omega_a,\\
 p^0\in\{0,1\}\Rightarrow p^+=p^0,\quad n^0=0\Rightarrow n^+=0,\\
 \epsilon_{\rm raw}=
 {|\delta{\cal E}_1(X_1^0)-\delta{\cal E}_2(X_2^-)-\tfrac12\Delta L(I_s^2-I_{\rm DC}^2)|\over E_*}\leq0.005,\\
 \epsilon_{\rm rec}=
 {4\sum_{jk}C_jw_ke_{jk}|p^+_{jk}-p^0_{jk}|
   +\sum_{ja}C_jc_a\bar\Omega_a|n^+_{ja}-n^0_{ja}|\over E_*}\leq0.01 .
\end{gather}
{\small \(\lambda\): raíz escalar del balance completo, con capacidad y
actividades positivas. No ajuste de calor durante la evolución,
equilibrio forzado ni una segunda absorción del fotón.}
\origin{Proyección y cambio de referencia propios.
\href{Supuestos.pdf\#nameddest=S27}{S27--S28};
\href{Parametros.pdf\#nameddest=P45}{P45--P47}.}
\clearpage
\block{Tiempo oficial y réplica 1D independiente}
\begin{gather}
 t_c=\text{primer cruce confirmado de }V_{\rm out}=0.1\ {\rm mV},\quad
 t_s=t_c+2\ {\rm ps},\quad \delta t_{\rm conf}=0.5\ {\rm ps},\\
 X^{\rm sh}_1(t_a)={\cal P}X_2(t_a),\quad 0\leq t_a-t_c\leq h_2,\quad h_2\leq0.02\ {\rm ps},\\
 \dot X^{\rm sh}_1=F_1(X^{\rm sh}_1),\quad t_a\leq t\leq t_s,\quad
 X_1(t_s^+)={\cal P}[X_2(t_s^-),X_1^{\rm sh}(t_s^-)],\\
 X_{\rm oficial}(t)=
 \begin{cases}X_2(t),&0\leq t\leq t_s,\\X_1(t),&t_s<t\leq15\ {\rm ns}.\end{cases}
\end{gather}
{\small Se reproyecta el último núcleo 2D y se conservan las colas nativas ya evolucionadas de la réplica. No se reinicializan frías al relevo. La fase global se alinea entre gauges con cortes derechos distintos; sus gradientes relativos no se eliminan.}
\block{Control longitudinal y errores del solapamiento}
\begin{gather}
 E_*(t)=E_{\rm ret}+\int_0^t|I_sV_{\rm dev}|\,{\rm d}t',\\
 \epsilon_{\rm cut}={\int_0^t(|P_L|+|P_R|)\,{\rm d}t'\over E_*(t)}\leq0.02,\quad
 P_\ell=\sum_{i\in\ell}\left(4C_i\sum_kw_ke_{ik}\dot p^{\rm raw}_{ik}
                    +\dot U^{\rm raw}_{{\rm ph},i}\right),\\
 \epsilon_d^2={\min_\theta\sum_iV_i|d_i^{(2)}-e^{i\theta}({\cal L}d_1)_i|^2
                    \over\sum_iV_i|d_i^{{\rm DC},2}|^2}\leq0.03^2,\\
 \epsilon_{\rm pop}={\|p_2-{\cal L}p_1,n_2-{\cal L}n_1\|_{E,1}\over E_*}\leq0.05,\\
 \epsilon_J^2={\sum_e (R_\square/c_e)(I_e-I_e^{1D})^2\over I_{\rm DC}^2R_n}
 \leq0.03^2,\quad
 I_e^{1D}={I_s\over W}L_e^{dual}{\Delta x_e\over\ell_e},\quad R_n=R_\square L/W,\\
 |V^{(2)}_{\rm out}-V^{\rm sh}_{\rm out}|\leq2\ \mu{\rm V},\quad
 {\max(|I_b^{(2)}-I_b^{\rm sh}|,|I_s^{(2)}-I_s^{\rm sh}|)\over I_{\rm DC}}\leq0.005,\\
 {|\delta{\cal E}_2+\tfrac12\Delta L(I_s^2-I_{\rm DC}^2)-\delta{\cal E}^{\rm sh}_1|\over E_*}\leq0.01,\quad
 N_+=N_-=0\quad(\text{contornos: }|d|>10^{-3}|d_{\rm DC}|).
\end{gather}
{\small \({\cal L}\): interpolación del campo relativo a DC y poblaciones.
Mezcla error transversal y de remapeo. \(P_L,P_R\) se calculan antes de
congelar reservorios. El flujo cortado indica truncamiento; no acota
matemáticamente el estado exterior desconocido.}
\block{Integración 1D y balance independiente}
\begin{gather}
 X_h=\Phi_h(X),\quad X_{h/2,h/2}=\Phi_{h/2}^2(X),\quad
 e_t={\|X_{h/2,h/2}-X_h\|_{\rm esc}\over3}\leq1,\\
 |\Delta E-\Delta W_{\rm mid}|
 \leq2.5\,10^{-6}\left(E_{\rm ret}+\sum|\Delta W_{\rm mid}|\right),\\
 \max_t{|E(t)+E_{\rm ref,shift}-E(0)-W_{\rm ext}(t)|
  \over E_{\rm ret}+\sum|\Delta W_{\rm mid}|}\leq0.01.
\end{gather}
\origin{Punto medio implícito acoplado heredado y duplicación temporal.
Ante fallo se conserva el estado con diagnóstico.
\href{Supuestos.pdf\#nameddest=S29}{S29--S32}.}
\end{document}
